\section{6 июня 1926 г.}

\lp{20260606_184920179}

Москва, 6 июня 1926 г.

Дорогой мой и милый Коля, не получив ни от тебя, ни от Анны Михайловны никакого
ответа на мои письма, я начинаю думать, что ты их не получил совсем; это было
бы мне весьма неприятно, т.к. придется тогда еще раз писать то же самое. Письма
эти я писал по твоему зимнему адресу: Fontaine \& Yvette. К сожалению я не
знал что на адресе надо писать «par Palaiseau» , может быть поэтому мои письма
и не дошли. Письма были посланы незаказными. Это письмо я пошлю заказным по
адресу, данному мне Соней и очень прошу тебя, не беря дурного примера с меня,
ответить открыткой по возможности тотчас же. (Если тебе некогда, попроси
сделать это Анну Михайловну).

Ответь мне на следующие вопросы: \\
1) При посылке тебе партитуры моей симфонии, я послал одновременно таковую и
Сергею Васильевичу по твоему адресу (его адрес я не знаю). Напиши получил ли он
ее. \\
2) Во второй половине апреля я послал тебе следующие ноты: мой струнный квартет
(партитуру и голоса), легкие этюды соч. 32 (две тетради),


\lp{20260606_184931545}

3 импровизации для виолончели с фортепиано, три тетради русских народных песен
стр 29, и органную прелюдию и фугу Баха в моем переложении для органа. Получил
ли ты все это? \\
3) Получил ли ты мое письмо от 28 апреля, посланное тебе в ответ на твое
письмо? \\
4) Одновременно с этим письмом я послал письмо Сергею Васильевичу по твоему
адресу, причем в этом письме, между прочим, я прошу его согласия на посвящение
ему моих двух прелюдий и фуг для органа. Они уже сданы в печать, а между тем
может быть это будет ему неприятно; и затем я очень прошу его не посылать мне
денег (он мне посылал почти каждый месяц по 10 долларов). Я его очень благодарю
за выплаты, но очень прошу этого не делать больше, т.к. мне предоставляется
возможным принимать помощь хотя бы и от близких людей лишь в случае болезни, а
пока я в силах работать, мне это тяжело и неприятно.
5) После смерти Егора Львовича Катуара я написал тебе длинное письмо, которое
может дойти до тебя лишь в том случае, если ты распорядился

\lp{20260606_184931545}

\noindent
о пересылке тебе писем по теперешнему твоему адресу, т.к. я не знал, что ты
уезжаешь, послал тебе это письмо в «Fontaine \& Yvette». В этом письме я
подробно изложил тебе обстоятельства, в которых оказались дочери Егора
Львовича, а также мои предложения и пожелания что можно было бы для них
сделать, и затем, в том же письме я прошу тебя и Сергея Васильевича, и Льва
Эдуардовича сделать все что от вас зависит для ознакомления заграничной публики
и музыкантов с сочинениями Егора Львовича (последние его сочинения: романсы на
тексты Соловьева, фортепианный квартет, фортепианный квинтет; Курс гармонии и в
рукописи большой труд по музык. форм., написанный карандашом. Кроме того нашлось
некоторое количество фортепианных пьес, совершенно законченных и ненапечатанных.
Не напечатаны его симфония, «Мцыри», и партитура «Русалки». В конце этого
письма, в связи с объявлением в ближайшем будущем конкурса на место профессора
по творческому отделению, я писал тебе, что единственным кандидатом, на котором
сошлись бы все любящие настоящую музыку, был бы ты и думается мне, что

\lp{20260606_184941420}

\noindent
хотя обстоятельства здесь довольно трудные (я имею ввиду не квартирные вопросы,
а главным образом современные «музыкальные уклоны»), все же в виде опыта может
ты и рискнул бы все же вернуться в Москву на несколько месяцев раньше, чем ты
предполагал (мне сказал Б. Б. Красин что ты с ним подписал договор на февраль
27-го года). Ответь мне, милый Коля, хотя бы в самых лаконических формах на это
письмо, и я тебе обещаю, что этим летом я не поленюсь отвечать тебе на твои
письма тотчас же, т.к. вопросы о руководителях учащихся творческого отдела
являются для меня вопросом судьбы всей музыкальной жизни нашей страны на многие
годы, и вопрос этот волнует меня не меньше, чем тебя. Сердечно кланяюсь Анне
Михайловне, Сергею Васильевичу, Нат. Алек., Льву Эдуардовичу, Ольге Николаевне.
Крепко тебя целую, \\ твой А. Гедике.

Екатерина Петровна и Паша также целуют тебя и \rem{...}.
Летом мы наверное никуда не поедем.

Пиши на конверте квартира № 15, а то письма неделями лежат в консерватории.
